% META: {"id": "1NSI-WEB-EX-022", "chapitre": "1NSI-WEB-IHM", "type_objet": "exercice", "capacites": ["1NSI-WEB-IHM-C4"], "parcours": 2, "competences": ["programmer", "analyser"], "duree_min": 8, "mode_creation": "ex_nihilo", "sources_inspiration": [], "fichier_tex": "chapitres/1NSI-WEB-IHM/exercices/1NSI-WEB-EX-022.tex", "corrige_tex": "chapitres/1NSI-WEB-IHM/corriges/1NSI-WEB-CO-022.tex", "status": "approved", "capacites_codes": ["C4"], "methodes": ["M1"]}
% BEGIN-VERIFY
% gestionnaires = {}
% def ajouter_gestionnaire(id_composant, fonction):
%     gestionnaires[id_composant] = fonction
% def declencher_clic(id_composant):
%     if id_composant in gestionnaires:
%         return gestionnaires[id_composant]()
%     return None
% etat = {"visible": False}
% def basculer_affichage():
%     etat["visible"] = not etat["visible"]
%     return etat["visible"]
% ajouter_gestionnaire("bouton_menu", basculer_affichage)
% resultats = [declencher_clic("bouton_menu") for _ in range(3)]
% assert resultats == [True, False, True]
% END-VERIFY
\begin{exercice}{1NSI-WEB-EX-002}{1}{12}
On reprend le modèle du cours qui enregistre un gestionnaire puis déclenche un clic.
Un menu de navigation doit s'afficher ou se masquer à chaque clic sur le bouton
\lstinline{"bouton_menu"} (un menu \og bascule \fg{}, comme sur un site mobile).
\begin{enumerate}
  \item Écrire une fonction \lstinline{basculer_affichage()} qui inverse la valeur de \lstinline{etat["visible"]} (initialement \lstinline{False}) et renvoie la nouvelle valeur.
  \item Associer cette fonction à \lstinline{"bouton_menu"} avec \lstinline{ajouter_gestionnaire}.
  \item Simuler $3$ clics successifs sur \lstinline{"bouton_menu"}. Quelle est la séquence de valeurs renvoyées ?
\end{enumerate}
\end{exercice}
